\label{sec:datos}

Los análisis propuestos en los distintos módulos se apoyan en un conjunto de fuentes 
de datos disponibles dentro de YPF. El siguiente cuadro resume qué bases son 
necesarias para cada módulo.

\begin{table}[h]
\centering
\caption{Fuentes de datos por módulo}
\label{tab:datos}
\setlength{\extrarowheight}{4pt}
\small
\begin{tabular}{lcccc}
\hline\hline
\textbf{Base de datos} & \textbf{Full} & \textbf{Serviclub} & 
\textbf{Cartelería} & \textbf{Boxes} \\
\hline
Dimensional estaciones / tiendas   & \checkmark & \checkmark & \checkmark & \checkmark \\
Transaccional estaciones            & \checkmark & \checkmark & \checkmark & \\
Transaccional tiendas               & \checkmark & \checkmark & \checkmark & \\
Transaccional Serviclub             &            & \checkmark & \checkmark & \\
Cartelería digital (contenidos)     &            &            & \checkmark & \\
Registros operativos boxes          &            &            &            & \checkmark \\
Datos de competidores / entorno     & \checkmark &            &            & \checkmark \\
Costos de inversión y operación     & \checkmark &            &            & \\
\hline\hline
\end{tabular}
\end{table}

A continuación se describe brevemente la información con la que se espera contar de cada fuente:

\begin{itemize}

    \item \textbf{Dimensional estaciones / tiendas.} Datos de ubicación y 
    características de cada estación: provincia, localidad, coordenadas 
    geográficas, nombre, tipo de red (propia, terceros, ACA), formato 
    (Full o tradicional), y variables demográficas del entorno cuando 
    estén disponibles (nivel socioeconómico, densidad poblacional).

    \item \textbf{Transaccional estaciones de servicio.} Datos a nivel de 
    despacho: tipo de combustible, litros y monto, fecha y hora, estación 
    y surtidor.

    \item \textbf{Transaccional tiendas.} Datos a nivel de transacción en 
    tienda: producto, categoría, monto, fecha y hora, estación.

    \item \textbf{Transaccional Serviclub.} Datos a nivel de transacción con 
    identificación de cliente: combustible, litros y monto, fecha y hora, 
    medio de pago (app, Mercado Pago, MODO u otro). Para compras en tienda 
    Full a través de la aplicación el detalle de producto.

    \item \textbf{Cartelería digital.} Registro de contenidos mostrados por 
    estación y franja horaria: tipo de mensaje, producto o categoría 
    promocionada, duración de la exposición. Esta base deberá construirse 
    o sistematizarse junto con el equipo responsable de la comunicación 
    comercial.

    \item \textbf{Registros operativos de boxes.} Información sobre capacidad 
    instalada y grado de utilización por estación y período, a partir de 
    los sistemas internos de gestión del servicio.

    \item \textbf{Datos de competidores y entorno.} Información 
    georreferenciada sobre estaciones competidoras y lubricentros 
    independientes, obtenida a partir de fuentes públicas, datos 
    administrativos o estimaciones de mercado disponibles.

    \item \textbf{Costos de inversión y operación.} Costos de remodelación 
    al formato Full y costos operativos adicionales por estación, necesarios 
    para el análisis de retorno de inversión del módulo~2.

\end{itemize}